\section{Method}
\label{sec:method}

InkFlow fine-tunes an open-source I2V backbone (Wan2.1-I2V) with
dual-stage conditioning and art-specific training objectives.
Figure~\ref{fig:pipeline} illustrates the overall framework.

\begin{figure}[t]
\centering
\fbox{\parbox{0.95\columnwidth}{\small
\textbf{[Figure placeholder: InkFlow pipeline diagram. Left: input
painting + user-provided line map and ink-density map. Center:
Wan2.1-I2V denoising process with stage-separated conditioning
(line map at steps $0$--$T/2$, ink-density map at steps
$T/2$--$T$). Right: generated video with blank-space preservation
and breathing motion.]}}}
\caption{InkFlow pipeline. The line-distribution map governs
structural composition during early denoising; the ink-density
distribution map governs tonal rendering during later denoising.
The blank-space loss prevents texture generation in
user-designated negative-space regions.}
\label{fig:pipeline}
\end{figure}

\subsection{Preliminaries: I2V Diffusion}

Given an input image $I_0$ and a text prompt $c$, an I2V diffusion
model generates a video $\{x_0, x_1, \ldots, x_L\}$ by iteratively
denoising from Gaussian noise. At each denoising step $t$, the
model predicts $\hat{x}_0$ from the noisy latent $x_t$:
\begin{equation}
\hat{x}_0 = f_\theta(x_t, t, I_0, c),
\end{equation}
where $f_\theta$ is the denoising network. We adopt Wan2.1-I2V as
the backbone and fine-tune it with additional conditioning and
loss terms.

\subsection{Dual-Stage Conditioning}
\label{sec:dual_condition}

\subsubsection{Line-Distribution Map}

The user provides a \emph{line-distribution map} $C_{\text{line}}$:
a binary or soft mask indicating the structural lines of the
composition (mountains, rivers, pavilions). This map is encoded
via a ControlNet-style adapter
$E_{\text{line}}$ and injected as additive conditioning during the
first half of denoising steps ($t \in [T/2, T]$):
\begin{equation}
\hat{x}_0 = f_\theta(x_t, t, I_0, c + E_{\text{line}}(C_{\text{line}})),
\quad t \geq T/2.
\end{equation}
This ensures that the structural composition is established before
tonal detail is rendered, mirroring the physical ink-wash workflow
where the artist first lays down the structural skeleton
(勾皴) before applying ink washes (渲染).

\subsubsection{Ink-Density Distribution Map}

The user provides an \emph{ink-density distribution map}
$C_{\text{ink}}$: a 5-channel map where each channel corresponds to
one of the five ink levels
(焦、浓、重、淡、清), with pixel values indicating the local density
of that ink level. This map is encoded via a separate adapter
$E_{\text{ink}}$ and injected during the second half of denoising
steps ($t \in [0, T/2)$):
\begin{equation}
\hat{x}_0 = f_\theta(x_t, t, I_0, c + E_{\text{ink}}(C_{\text{ink}})),
\quad t < T/2.
\end{equation}

\subsubsection{Five-Level Ink Histogram Matching}

For precise control mode, we enforce that the generated video
frames match a target 5-bin ink-density histogram $H^*$ derived
from the ink-density map. At each training step, we compute the
histogram $H(\hat{x}_0)$ of the predicted frame and minimize:
\begin{equation}
\mathcal{L}_{\text{ink}} = \text{EMD}(H(\hat{x}_0), H^*),
\end{equation}
where EMD is the Earth Mover's Distance between the predicted and
target ink-level histograms. This provides a differentiable
measure of ink-density fidelity.

\subsection{Blank-Space Loss}
\label{sec:blank_space}

We formalize the aesthetic principle of \emph{treating white as
black} (计白当黑) as a training objective. The user designates
negative-space regions via a binary mask $M_{\text{blank}}$. The
blank-space loss penalizes any generated texture inside these
regions:
\begin{equation}
\mathcal{L}_{\text{blank}} = \frac{1}{|M_{\text{blank}}|}
\sum_{(x,y) \in M_{\text{blank}}} \text{Var}_{\text{local}}(\hat{x}_0(x,y)),
\end{equation}
where $\text{Var}_{\text{local}}$ measures the local pixel-variance
in a $5 \times 5$ window. This encourages the model to keep
negative-space regions flat and texture-free, preserving the
compositional voids that are essential to Shanshui aesthetics.

\subsection{Rhythm-of-Strokes Temporal Prior}
\label{sec:rhythm}

Ink-wash animation is characterized by gentle, breathing-like
motion---drifting mist, slowly flowing water---rather than abrupt
movement. We encourage this quality via a temporal smoothness
prior inspired by calligraphic timing:
\begin{equation}
\mathcal{L}_{\text{rhythm}} = \sum_{l=1}^{L-1}
\| x_l - x_{l-1} - \mu_{\text{flow}} \|^2,
\end{equation}
where $\mu_{\text{flow}}$ is a slowly varying optical-flow
estimate between consecutive frames. This penalizes high-frequency
temporal jitter while allowing smooth, low-frequency motion.

\subsection{Training Objective}

The total training objective combines the standard I2V diffusion
loss with our art-specific terms:
\begin{equation}
\mathcal{L} = \mathcal{L}_{\text{diff}} + \lambda_{\text{ink}}
\mathcal{L}_{\text{ink}} + \lambda_{\text{blank}}
\mathcal{L}_{\text{blank}} + \lambda_{\text{rhythm}}
\mathcal{L}_{\text{rhythm}},
\end{equation}
where $\lambda_{\text{ink}} = 0.5$, $\lambda_{\text{blank}} = 0.3$,
and $\lambda_{\text{rhythm}} = 0.2$ by default. The model is
fine-tuned on a dataset of 30 Shanshui paintings with
user-annotated line maps, ink-density maps, and blank-space masks.

\subsection{Implementation Details}

We build on the existing OilPaintingVideo project's ComfyUI
pipeline for the base I2V infrastructure. The dual-condition
adapters are implemented as lightweight ControlNet branches
(${\sim}5$\,M parameters each) attached to the Wan2.1-I2V U-Net.
Training uses AdamW with learning rate $1 \times 10^{-5}$, batch
size 1, and 4-frame video clips on a single NVIDIA RTX A4500 GPU.
The blank-space mask and ink-density map are extracted
semi-automatically: a U-Net segmenter identifies candidate
blank-space regions, which the user can refine; the ink-density
map is computed from the input painting's luminance histogram
quantized into 5 bins aligned to the five ink levels.
